\documentclass[conference]{IEEEtran}
\usepackage{cite}
\usepackage{amsmath,amssymb,amsfonts}
\usepackage{graphicx}
\usepackage{textcomp}
\usepackage{xcolor}
% Project additions to the IEEE template.
\usepackage[T1]{fontenc}
\usepackage{booktabs}
\usepackage{url}
\usepackage{tikz}
\usetikzlibrary{positioning,arrows.meta}
\usepackage{pgfplots}
\usepgfplotslibrary{groupplots}
\pgfplotsset{compat=1.18}
\newcommand{\NaR}{\mathrm{NaR}}
\newcommand{\code}[1]{\texttt{#1}}
\makeatletter
\newcommand{\linebreakand}{\end{@IEEEauthorhalign}\hfill\mbox{}\par
  \mbox{}\hfill\begin{@IEEEauthorhalign}}
\makeatother
\usepackage[hidelinks]{hyperref}
\usepackage{etoolbox}
\patchcmd{\abstract}{---}{: }{}{\errmessage{Abstract punctuation patch failed}}
\patchcmd{\IEEEkeywords}{---}{: }{}{\errmessage{Keywords punctuation patch failed}}
% Benchmark numbers come from scripts/summarize_benchmark.py. Until the
% real benchmark run exists, the synthetic offline run is used and marked.
\IfFileExists{../results/spmv-macros.tex}{%
  % Generated by scripts/summarize_benchmark.py from spmv-errors.csv
\newcommand{\BenchMatrices}{358}
\newcommand{\BenchNnz}{3{,}054{,}096}
\newcommand{\BenchMaxNnz}{49{,}920}
\newcommand{\BenchMedianTEight}{-0.7}
\newcommand{\BenchConvMedianTEight}{-1.2}
\newcommand{\BenchBelowOneTEight}{87}
\newcommand{\BenchInfTEight}{0}
\newcommand{\BenchSatTEight}{0}
\newcommand{\BenchOverflowTEight}{0}
\newcommand{\BenchFlushTEight}{0}
\newcommand{\BenchAboveOneTEight}{13}
\newcommand{\BenchVecPerNnzTEight}{1.23}
\newcommand{\BenchMedianEFourMThree}{-1.2}
\newcommand{\BenchConvMedianEFourMThree}{-1.6}
\newcommand{\BenchBelowOneEFourMThree}{64}
\newcommand{\BenchInfEFourMThree}{34}
\newcommand{\BenchSatEFourMThree}{0}
\newcommand{\BenchOverflowEFourMThree}{120}
\newcommand{\BenchFlushEFourMThree}{173}
\newcommand{\BenchAboveOneEFourMThree}{36}
\newcommand{\BenchVecPerNnzEFourMThree}{1.80}
\newcommand{\BenchVecReductionEFourMThree}{32}
\newcommand{\BenchTotalReductionEFourMThree}{16}
\newcommand{\BenchMedianEFiveMTwo}{-1.0}
\newcommand{\BenchConvMedianEFiveMTwo}{-1.3}
\newcommand{\BenchBelowOneEFiveMTwo}{73}
\newcommand{\BenchInfEFiveMTwo}{25}
\newcommand{\BenchSatEFiveMTwo}{0}
\newcommand{\BenchOverflowEFiveMTwo}{86}
\newcommand{\BenchFlushEFiveMTwo}{138}
\newcommand{\BenchAboveOneEFiveMTwo}{27}
\newcommand{\BenchVecPerNnzEFiveMTwo}{1.80}
\newcommand{\BenchVecReductionEFiveMTwo}{32}
\newcommand{\BenchTotalReductionEFiveMTwo}{16}
\newcommand{\BenchMedianTSixteen}{-3.2}
\newcommand{\BenchConvMedianTSixteen}{-3.6}
\newcommand{\BenchBelowOneTSixteen}{97}
\newcommand{\BenchInfTSixteen}{0}
\newcommand{\BenchSatTSixteen}{0}
\newcommand{\BenchOverflowTSixteen}{0}
\newcommand{\BenchFlushTSixteen}{0}
\newcommand{\BenchAboveOneTSixteen}{3}
\newcommand{\BenchVecPerNnzTSixteen}{1.47}
\newcommand{\BenchMedianFPSixteen}{-3.5}
\newcommand{\BenchConvMedianFPSixteen}{-3.7}
\newcommand{\BenchBelowOneFPSixteen}{75}
\newcommand{\BenchInfFPSixteen}{25}
\newcommand{\BenchSatFPSixteen}{0}
\newcommand{\BenchOverflowFPSixteen}{86}
\newcommand{\BenchFlushFPSixteen}{102}
\newcommand{\BenchAboveOneFPSixteen}{25}
\newcommand{\BenchVecPerNnzFPSixteen}{1.84}
\newcommand{\BenchVecReductionFPSixteen}{20}
\newcommand{\BenchTotalReductionFPSixteen}{9}
\newcommand{\BenchMedianTThirtyTwo}{-8.2}
\newcommand{\BenchConvMedianTThirtyTwo}{-8.6}
\newcommand{\BenchBelowOneTThirtyTwo}{100}
\newcommand{\BenchInfTThirtyTwo}{0}
\newcommand{\BenchSatTThirtyTwo}{0}
\newcommand{\BenchOverflowTThirtyTwo}{0}
\newcommand{\BenchFlushTThirtyTwo}{0}
\newcommand{\BenchAboveOneTThirtyTwo}{0}
\newcommand{\BenchVecPerNnzTThirtyTwo}{2.14}
\newcommand{\BenchMedianFPThirtyTwo}{-7.6}
\newcommand{\BenchConvMedianFPThirtyTwo}{-7.8}
\newcommand{\BenchBelowOneFPThirtyTwo}{100}
\newcommand{\BenchInfFPThirtyTwo}{0}
\newcommand{\BenchSatFPThirtyTwo}{0}
\newcommand{\BenchOverflowFPThirtyTwo}{0}
\newcommand{\BenchFlushFPThirtyTwo}{2}
\newcommand{\BenchAboveOneFPThirtyTwo}{0}
\newcommand{\BenchVecPerNnzFPThirtyTwo}{2.76}
\newcommand{\BenchVecReductionFPThirtyTwo}{22}
\newcommand{\BenchTotalReductionFPThirtyTwo}{11}
\newcommand{\BenchMedianTEightStep}{-0.7}
\newcommand{\BenchConvMedianTEightStep}{-1.2}
\newcommand{\BenchBelowOneTEightStep}{87}
\newcommand{\BenchInfTEightStep}{0}
\newcommand{\BenchSatTEightStep}{0}
\newcommand{\BenchOverflowTEightStep}{0}
\newcommand{\BenchFlushTEightStep}{0}
\newcommand{\BenchAboveOneTEightStep}{13}
\newcommand{\BenchVecPerNnzTEightStep}{1.23}
\newcommand{\BenchMedianTSixteenStep}{-3.2}
\newcommand{\BenchConvMedianTSixteenStep}{-3.6}
\newcommand{\BenchBelowOneTSixteenStep}{97}
\newcommand{\BenchInfTSixteenStep}{0}
\newcommand{\BenchSatTSixteenStep}{0}
\newcommand{\BenchOverflowTSixteenStep}{0}
\newcommand{\BenchFlushTSixteenStep}{0}
\newcommand{\BenchAboveOneTSixteenStep}{3}
\newcommand{\BenchVecPerNnzTSixteenStep}{1.47}
\newcommand{\BenchMedianTThirtyTwoStep}{-8.2}
\newcommand{\BenchConvMedianTThirtyTwoStep}{-8.6}
\newcommand{\BenchBelowOneTThirtyTwoStep}{100}
\newcommand{\BenchInfTThirtyTwoStep}{0}
\newcommand{\BenchSatTThirtyTwoStep}{0}
\newcommand{\BenchOverflowTThirtyTwoStep}{0}
\newcommand{\BenchFlushTThirtyTwoStep}{0}
\newcommand{\BenchAboveOneTThirtyTwoStep}{0}
\newcommand{\BenchVecPerNnzTThirtyTwoStep}{2.14}
\newcommand{\BenchInRangeTEightEFourMThree}{235}
\newcommand{\BenchBetterTEightEFourMThree}{13}
\newcommand{\BenchWorseTEightEFourMThree}{181}
\newcommand{\BenchEqualTEightEFourMThree}{41}
\newcommand{\BenchInRangeTEightEFiveMTwo}{269}
\newcommand{\BenchBetterTEightEFiveMTwo}{45}
\newcommand{\BenchWorseTEightEFiveMTwo}{167}
\newcommand{\BenchEqualTEightEFiveMTwo}{57}
\newcommand{\BenchInRangeTSixteenFPSixteen}{270}
\newcommand{\BenchBetterTSixteenFPSixteen}{48}
\newcommand{\BenchWorseTSixteenFPSixteen}{110}
\newcommand{\BenchEqualTSixteenFPSixteen}{112}
\newcommand{\BenchInRangeTThirtyTwoFPThirtyTwo}{358}
\newcommand{\BenchBetterTThirtyTwoFPThirtyTwo}{199}
\newcommand{\BenchWorseTThirtyTwoFPThirtyTwo}{33}
\newcommand{\BenchEqualTThirtyTwoFPThirtyTwo}{126}
\newcommand{\BenchInRangeTEightTEightStep}{358}
\newcommand{\BenchBetterTEightTEightStep}{52}
\newcommand{\BenchWorseTEightTEightStep}{51}
\newcommand{\BenchEqualTEightTEightStep}{255}
\newcommand{\BenchInRangeTSixteenTSixteenStep}{358}
\newcommand{\BenchBetterTSixteenTSixteenStep}{10}
\newcommand{\BenchWorseTSixteenTSixteenStep}{13}
\newcommand{\BenchEqualTSixteenTSixteenStep}{335}
\newcommand{\BenchInRangeTThirtyTwoTThirtyTwoStep}{358}
\newcommand{\BenchBetterTThirtyTwoTThirtyTwoStep}{5}
\newcommand{\BenchWorseTThirtyTwoTThirtyTwoStep}{0}
\newcommand{\BenchEqualTThirtyTwoTThirtyTwoStep}{353}
\def\BenchCDF{../results/spmv-cdf.csv}\def\BenchDraft{}}{%
  % Generated by scripts/summarize_benchmark.py from synthetic-errors.csv
\newcommand{\BenchMatrices}{64}
\newcommand{\BenchNnz}{122{,}151}
\newcommand{\BenchMaxNnz}{6{,}712}
\newcommand{\BenchMedianTEight}{-0.6}
\newcommand{\BenchConvMedianTEight}{-0.8}
\newcommand{\BenchBelowOneTEight}{100}
\newcommand{\BenchInfTEight}{0}
\newcommand{\BenchSatTEight}{0}
\newcommand{\BenchOverflowTEight}{0}
\newcommand{\BenchFlushTEight}{0}
\newcommand{\BenchAboveOneTEight}{0}
\newcommand{\BenchVecPerNnzTEight}{0.98}
\newcommand{\BenchMedianEFourMThree}{\infty}
\newcommand{\BenchConvMedianEFourMThree}{-1.6}
\newcommand{\BenchBelowOneEFourMThree}{39}
\newcommand{\BenchInfEFourMThree}{61}
\newcommand{\BenchSatEFourMThree}{0}
\newcommand{\BenchOverflowEFourMThree}{30}
\newcommand{\BenchFlushEFourMThree}{57}
\newcommand{\BenchAboveOneEFourMThree}{61}
\newcommand{\BenchVecPerNnzEFourMThree}{1.42}
\newcommand{\BenchVecReductionEFourMThree}{32}
\newcommand{\BenchTotalReductionEFourMThree}{16}
\newcommand{\BenchMedianEFiveMTwo}{-1.1}
\newcommand{\BenchConvMedianEFiveMTwo}{-1.3}
\newcommand{\BenchBelowOneEFiveMTwo}{55}
\newcommand{\BenchInfEFiveMTwo}{45}
\newcommand{\BenchSatEFiveMTwo}{0}
\newcommand{\BenchOverflowEFiveMTwo}{28}
\newcommand{\BenchFlushEFiveMTwo}{35}
\newcommand{\BenchAboveOneEFiveMTwo}{45}
\newcommand{\BenchVecPerNnzEFiveMTwo}{1.42}
\newcommand{\BenchVecReductionEFiveMTwo}{32}
\newcommand{\BenchTotalReductionEFiveMTwo}{16}
\newcommand{\BenchMedianTSixteen}{-2.9}
\newcommand{\BenchConvMedianTSixteen}{-3.2}
\newcommand{\BenchBelowOneTSixteen}{100}
\newcommand{\BenchInfTSixteen}{0}
\newcommand{\BenchSatTSixteen}{0}
\newcommand{\BenchOverflowTSixteen}{0}
\newcommand{\BenchFlushTSixteen}{0}
\newcommand{\BenchAboveOneTSixteen}{0}
\newcommand{\BenchVecPerNnzTSixteen}{1.28}
\newcommand{\BenchMedianFPSixteen}{-3.5}
\newcommand{\BenchConvMedianFPSixteen}{-3.7}
\newcommand{\BenchBelowOneFPSixteen}{56}
\newcommand{\BenchInfFPSixteen}{44}
\newcommand{\BenchSatFPSixteen}{0}
\newcommand{\BenchOverflowFPSixteen}{19}
\newcommand{\BenchFlushFPSixteen}{21}
\newcommand{\BenchAboveOneFPSixteen}{44}
\newcommand{\BenchVecPerNnzFPSixteen}{1.62}
\newcommand{\BenchVecReductionFPSixteen}{21}
\newcommand{\BenchTotalReductionFPSixteen}{10}
\newcommand{\BenchMedianTThirtyTwo}{-7.8}
\newcommand{\BenchConvMedianTThirtyTwo}{-8.0}
\newcommand{\BenchBelowOneTThirtyTwo}{100}
\newcommand{\BenchInfTThirtyTwo}{0}
\newcommand{\BenchSatTThirtyTwo}{0}
\newcommand{\BenchOverflowTThirtyTwo}{0}
\newcommand{\BenchFlushTThirtyTwo}{0}
\newcommand{\BenchAboveOneTThirtyTwo}{0}
\newcommand{\BenchVecPerNnzTThirtyTwo}{1.92}
\newcommand{\BenchMedianFPThirtyTwo}{-7.5}
\newcommand{\BenchConvMedianFPThirtyTwo}{-7.6}
\newcommand{\BenchBelowOneFPThirtyTwo}{100}
\newcommand{\BenchInfFPThirtyTwo}{0}
\newcommand{\BenchSatFPThirtyTwo}{0}
\newcommand{\BenchOverflowFPThirtyTwo}{0}
\newcommand{\BenchFlushFPThirtyTwo}{0}
\newcommand{\BenchAboveOneFPThirtyTwo}{0}
\newcommand{\BenchVecPerNnzFPThirtyTwo}{2.49}
\newcommand{\BenchVecReductionFPThirtyTwo}{23}
\newcommand{\BenchTotalReductionFPThirtyTwo}{11}
\newcommand{\BenchMedianTEightStep}{-0.6}
\newcommand{\BenchConvMedianTEightStep}{-0.8}
\newcommand{\BenchBelowOneTEightStep}{100}
\newcommand{\BenchInfTEightStep}{0}
\newcommand{\BenchSatTEightStep}{0}
\newcommand{\BenchOverflowTEightStep}{0}
\newcommand{\BenchFlushTEightStep}{0}
\newcommand{\BenchAboveOneTEightStep}{0}
\newcommand{\BenchVecPerNnzTEightStep}{0.98}
\newcommand{\BenchMedianTSixteenStep}{-2.9}
\newcommand{\BenchConvMedianTSixteenStep}{-3.2}
\newcommand{\BenchBelowOneTSixteenStep}{100}
\newcommand{\BenchInfTSixteenStep}{0}
\newcommand{\BenchSatTSixteenStep}{0}
\newcommand{\BenchOverflowTSixteenStep}{0}
\newcommand{\BenchFlushTSixteenStep}{0}
\newcommand{\BenchAboveOneTSixteenStep}{0}
\newcommand{\BenchVecPerNnzTSixteenStep}{1.28}
\newcommand{\BenchMedianTThirtyTwoStep}{-7.8}
\newcommand{\BenchConvMedianTThirtyTwoStep}{-8.0}
\newcommand{\BenchBelowOneTThirtyTwoStep}{100}
\newcommand{\BenchInfTThirtyTwoStep}{0}
\newcommand{\BenchSatTThirtyTwoStep}{0}
\newcommand{\BenchOverflowTThirtyTwoStep}{0}
\newcommand{\BenchFlushTThirtyTwoStep}{0}
\newcommand{\BenchAboveOneTThirtyTwoStep}{0}
\newcommand{\BenchVecPerNnzTThirtyTwoStep}{1.92}
\newcommand{\BenchInRangeTEightEFourMThree}{25}
\newcommand{\BenchBetterTEightEFourMThree}{0}
\newcommand{\BenchWorseTEightEFourMThree}{25}
\newcommand{\BenchEqualTEightEFourMThree}{0}
\newcommand{\BenchInRangeTEightEFiveMTwo}{35}
\newcommand{\BenchBetterTEightEFiveMTwo}{9}
\newcommand{\BenchWorseTEightEFiveMTwo}{26}
\newcommand{\BenchEqualTEightEFiveMTwo}{0}
\newcommand{\BenchInRangeTSixteenFPSixteen}{36}
\newcommand{\BenchBetterTSixteenFPSixteen}{6}
\newcommand{\BenchWorseTSixteenFPSixteen}{30}
\newcommand{\BenchEqualTSixteenFPSixteen}{0}
\newcommand{\BenchInRangeTThirtyTwoFPThirtyTwo}{64}
\newcommand{\BenchBetterTThirtyTwoFPThirtyTwo}{45}
\newcommand{\BenchWorseTThirtyTwoFPThirtyTwo}{19}
\newcommand{\BenchEqualTThirtyTwoFPThirtyTwo}{0}
\newcommand{\BenchInRangeTEightTEightStep}{64}
\newcommand{\BenchBetterTEightTEightStep}{19}
\newcommand{\BenchWorseTEightTEightStep}{17}
\newcommand{\BenchEqualTEightTEightStep}{28}
\newcommand{\BenchInRangeTSixteenTSixteenStep}{64}
\newcommand{\BenchBetterTSixteenTSixteenStep}{0}
\newcommand{\BenchWorseTSixteenTSixteenStep}{0}
\newcommand{\BenchEqualTSixteenTSixteenStep}{64}
\newcommand{\BenchInRangeTThirtyTwoTThirtyTwoStep}{64}
\newcommand{\BenchBetterTThirtyTwoTThirtyTwoStep}{0}
\newcommand{\BenchWorseTThirtyTwoTThirtyTwoStep}{0}
\newcommand{\BenchEqualTThirtyTwoTThirtyTwoStep}{64}
\def\BenchCDF{../results/synthetic-cdf.csv}%
  \def\BenchDraft{\textcolor{red}{[DRAFT: synthetic matrices]} }}
\hypersetup{pdftitle={Takum Arithmetic in the RISC-V Vector Extension},
  pdfauthor={A. S. Fernandes, B. L. Soares, L. R. P. Ayala, P. H. C. Lima, P. L. Ozelin}}
\begin{document}

\title{Takum Arithmetic in the RISC-V Vector Extension: Instruction Design and
Sparse Matrix-Vector Evaluation}
\author{\IEEEauthorblockN{Arthur dos Santos Fernandes}
\IEEEauthorblockA{\textit{Department of Electronic Engineering} \\
\textit{Universidade Federal de Minas Gerais}\\
Belo Horizonte, Brazil}
\and
\IEEEauthorblockN{Bruno Lima Soares}
\IEEEauthorblockA{\textit{Department of Electronic Engineering} \\
\textit{Universidade Federal de Minas Gerais}\\
Belo Horizonte, Brazil}
\and
\IEEEauthorblockN{Lucas Rodrigues Prosperi Ayala}
\IEEEauthorblockA{\textit{Department of Electronic Engineering} \\
\textit{Universidade Federal de Minas Gerais}\\
Belo Horizonte, Brazil}
\linebreakand
\IEEEauthorblockN{Pedro Henrique Campos Lima}
\IEEEauthorblockA{\textit{Department of Electronic Engineering} \\
\textit{Universidade Federal de Minas Gerais}\\
Belo Horizonte, Brazil}
\and
\IEEEauthorblockN{Pedro Lopes Ozelin}
\IEEEauthorblockA{\textit{Department of Electronic Engineering} \\
\textit{Universidade Federal de Minas Gerais}\\
Belo Horizonte, Brazil}
}

\maketitle

\begin{abstract}
\BenchDraft Takum is a tapered-precision number format with one decoding scheme
for every bit width. We study it as the element type of a simplified RISC-V
vector extension. An exact functional model implements linear takums at 8, 16,
32, and 64 bits and five proposed instructions: element-wise add, subtract, and
multiply, a dot product that accumulates exactly and rounds once to twice the
element width, and a narrowing conversion. The instructions reuse the RVV
encoding layout under a custom opcode and need no new control state: the element
width selects the format, and saturation and the single NaR value replace the
rounding mode and exception flags of IEEE 754. We compare a CSR sparse
matrix-vector kernel in takum and in OFP8, binary16, and binary32 over
\BenchMatrices{} real matrices from the Matrix Market, measuring the relative
error against a binary128 reference and the dynamic instruction count. Takum
results were always finite, while E4M3, E5M2, and binary16 produced a NaN or an
infinity for \BenchInfEFourMThree\%, \BenchInfEFiveMTwo\%, and
\BenchInfFPSixteen\% of the matrices. Where the IEEE formats stay in range they
are usually more accurate at 8 and 16 bits; at 32 bits takum is more accurate
for \BenchBetterTThirtyTwoFPThirtyTwo{} matrices and less for
\BenchWorseTThirtyTwoFPThirtyTwo. The fused dot product removes
\BenchVecReductionFPSixteen\% to \BenchVecReductionEFourMThree\% of the vector
instructions of the kernel.
\end{abstract}
\begin{IEEEkeywords}
Takum, tapered precision, RISC-V vector extension, SpMV, functional simulation
\end{IEEEkeywords}

\section{Introduction}
Vector units now support many low-precision floating-point formats: IEEE 754
binary16~\cite{ieee754}, bfloat16, and the OFP8 types E4M3 and
E5M2~\cite{ofp8}. Each keeps the IEEE 754 framework and its exceptional values,
and supporting all of them makes the instruction set larger and less
regular~\cite{simd}. Takum arithmetic~\cite{takum} offers one alternative
family: its dynamic range is almost independent of the bit width, and all widths
share a decoder that reads at most the 12 most significant bits. In a study of
Intel AVX10.2, Hunhold found takums more accurate than the AVX10.2 formats over
sparse matrices from the SuiteSparse collection and proposed a smaller
instruction set built on them, naming RISC-V vector extensions as future
work~\cite{simd}.

This paper takes up that suggestion for the RISC-V Vector extension
(RVV)~\cite{rvv}. The project had three milestones: an exact scalar model of
linear takums, a vector extension with five instructions (Zvtakum), and an
experimental comparison with IEEE formats on sparse matrix-vector
multiplication (SpMV). Our contributions are (i) an encoding and semantics for
the five instructions that adds no control state to the ISA, (ii) a functional
vector model in which the takum and IEEE kernels run under one instruction
accounting, and (iii) an evaluation of accuracy and instruction count over the
real sparse matrices. We make no claims about cycles,
area, or power.

\section{Related Work}
Posits~\cite{posit} introduced tapered precision as an alternative to IEEE 754.
Takums keep the accuracy of posits near unity while encoding numbers far from
it more efficiently~\cite{takum}. Linear takums interpret the last field as a
fraction~$F$; the logarithmic variant uses the same fields with a different
value function~\cite{takum}. We implement the linear variant, which the AVX10.2
study also uses~\cite{simd}.

Hunhold proposed fused takum dot products that accumulate in twice the input
width (\code{VDPPT8PT16} to \code{VDPPT32PT64}) and showed that one takum
family removes most of the format-specific conversion instructions of
AVX10.2~\cite{simd}. Takums have also been evaluated in sparse linear
solvers~\cite{sparsesolvers}, and a VHDL codec has been
described~\cite{takumvhdl}. In RISC-V, FPHUB~\cite{fphub} also departs from
strict IEEE 754 compliance to reduce hardware cost, keeping an IEEE-like layout
with a single rounding mode. The ratified RVV has no FP8 arithmetic; we
therefore model FP8 conversions explicitly (Section~\ref{sec:baseline}).

\section{Takum Representation and Scalar Arithmetic}
An $n$-bit takum holds a sign bit $S$, a direction bit $D$, three regime bits
read as $R$, $r$ characteristic bits read as $C$, and $p=n-5-r$ fraction bits
read as $F$. With
\begin{equation}
 r=\begin{cases}R,&D=1\\7-R,&D=0\end{cases},\quad
 c=\begin{cases}2^r-1+C,&D=1\\1-2^{r+1}+C,&D=0\end{cases},
\end{equation}
$e=(1-2S)(c+S)$ and $f=F2^{-p}$, a finite nonzero value is
$x=[(1-3S)+f]\,2^e$. Short strings are extended with zero \emph{ghost bits}, so
the same definition serves T8. All-zero bits encode $0$, a lone sign bit encodes
$\NaR$ (not a real), and negation is two's complement.

The scalar model (\code{TakumFormat}, \code{TakumALU}) computes with exact
rationals and rounds each result once, to nearest with ties to even on the bit
stream, as \code{libtakum}~\cite{libtakum} does. Finite values outside the range
saturate to the largest or smallest magnitude and never become $0$ or $\NaR$;
$\NaR$ propagates through every operation. The first milestone validated this
model with an independent oracle for all $256^2$ T8 operand pairs, exhaustive T8
and T16 round trips, and 548,793 differential comparisons against a pinned
\code{libtakum} revision, with no mismatch. T64, needed as the accumulator of
T32, uses the same width-parametric code; tests check that zero extension of
every T8 pattern and of sampled T16 and T32 patterns preserves the value.

\section{Proposed Architecture: Zvtakum}
\subsection{Vector state and instructions}
Zvtakum keeps the RVV 1.0 state unchanged: 32 vector registers of VLEN bits,
\code{vl} and \code{vtype}. As for IEEE instructions, the element width SEW
selects the format: SEW${}=8,16,32$ means T8, T16, T32. Table~\ref{tab:enc}
lists the five instructions. They use the OP-V field layout
(\code{funct6|vm|vs2|vs1|funct3|vd|opcode}) under the \emph{custom-2} major
opcode, which RISC-V reserves for nonstandard extensions, so they cannot clash
with standard encodings. Each \code{funct6} repeats that of the IEEE
counterpart, and \code{funct3} is OPFVV.

\begin{table}[t]
\caption{Zvtakum instructions (opcode custom-2, funct3 OPFVV). funct6 repeats
that of \code{vfadd}, \code{vfsub}, \code{vfmul}, \code{vfwmacc}, and VFUNARY0}
\label{tab:enc}
\centering
\setlength{\tabcolsep}{4pt}
\begin{tabular}{lll}
\toprule
Instruction & funct6 & Semantics\\
\midrule
\code{vtadd.vv} & \code{000000} & $v_d[i]=\mathrm{rd}(v_{s2}[i]+v_{s1}[i])$\\
\code{vtsub.vv} & \code{000010} & $v_d[i]=\mathrm{rd}(v_{s2}[i]-v_{s1}[i])$\\
\code{vtmul.vv} & \code{100100} & $v_d[i]=\mathrm{rd}(v_{s2}[i]\,v_{s1}[i])$\\
\code{vtdot.vv} & \code{111100} & $v_d[0]=\mathrm{rd}_2(v_d[0]+\textstyle\sum_i v_{s2}[i]\,v_{s1}[i])$\\
\code{vncvt.t.t} & \code{010010} & $v_d[i]=\mathrm{rd}(v_{s2}[i])$, $2\mathrm{SEW}\to\mathrm{SEW}$\\
\bottomrule
\end{tabular}
\end{table}

Element-wise instructions apply the scalar ALU to the active elements; masked
and tail elements are left undisturbed. \code{vtdot.vv} reads SEW-bit sources
and a $2\mathrm{SEW}$-bit accumulator in $v_d[0]$ (T8 to T16, T16 to T32, T32
to T64); the sum of products is exact and rounded once ($\mathrm{rd}_2$), as in
Hunhold's fused dot products~\cite{simd}. \code{vncvt.t.t} narrows with the
same rounding and saturation; widening is exact zero extension and needs no
instruction in our kernels.

\subsection{Control and status state}
The F extension couples every IEEE instruction to \code{fcsr}: a 3-bit
rounding mode \code{frm} with five modes and the flags NV, DZ, OF, UF, and
NX~\cite{rvunpriv}. Zvtakum reads and writes none of it. There is one rounding
rule. Invalid results are $\NaR$, carried in the data. Overflow and underflow
saturate, so results stay finite and ordered. There are no subnormals, signed
zeros, infinities, or NaN payloads. Selecting E4M3 or E5M2 at SEW${}=8$ needs
one more bit of format state in an IEEE design; with takums one format per width
suffices.

\subsection{Functional vector model}
\code{VectorMachine} holds the register file as bytes, checks register-group
alignment, and implements the RVV subset used by the kernels: \code{vsetvli},
unit-stride and indexed loads and stores, \code{vmv.s.x}, and the IEEE baseline.
Every instruction is counted, and scalar loop instructions are reported by the
kernel so that totals are complete. The model is functional: it fixes
architectural results, not timing.

\subsection{SpMV kernel and IEEE baseline}\label{sec:baseline}
The CSR kernel processes one row at a time. Per strip of at most VLMAX
nonzeros it loads the values and the column offsets, gathers $x$ with
\code{vluxei32}, and accumulates. The takum body is one \code{vtdot.vv}. The
binary16 and binary32 body is \code{vfwmul.vv}, a \code{vsetvli} to $2\mathrm{SEW}$,
and an ordered \code{vfredosum.vs}; for FP8 two \code{vfwcvt} widen the operands
to binary16 first, because RVV has no FP8 arithmetic. Every format thus
accumulates in $2\mathrm{SEW}$ (FP8 in binary16, binary16 in binary32, binary32
in binary64), and a final narrowing pass converts $y$ to SEW. Each strip costs
5 vector instructions with takums, 7 with binary16 or binary32, and 9 with FP8,
plus the same scalar loop control. E4M3 follows OFP8 in non-saturating mode, so
overflow gives NaN.

\section{Evaluation Methodology}
\emph{Matrices.} We planned to use, as~\cite{simd} did, every real matrix of
the SuiteSparse collection~\cite{suitesparse} with at most 50,000 nonzeros, but
its file server refused every download during the project. We therefore use the
NIST Matrix Market~\cite{matrixmarket}, whose Harwell-Boeing, SPARSKIT, NEP,
and miscellaneous sets are also part of SuiteSparse: every real, integer, or
pattern matrix in coordinate format with at most 50,000 nonzeros after
symmetric expansion, \BenchMatrices{} matrices with \BenchNnz{} nonzeros in
total. The set is smaller and older than that of~\cite{simd}. Values are read
from the Matrix Market text as exact decimals.

\emph{Reference and error.} With $x$ the vector of ones, the reference
$y_{\mathrm{ref}}$ rounds the inputs, every product, and every sum to binary128.
For each format we report the relative error
$\|y_f-y_{\mathrm{ref}}\|_2/\|y_{\mathrm{ref}}\|_2$ and, as in~\cite{simd}, the
conversion error $\|A_f-A\|_F/\|A\|_F$. Norms are computed exactly with
integers, so the measurement adds no error. A NaN or infinity in $y_f$ counts as
an infinite error. Matrices with $y_{\mathrm{ref}}=0$ are excluded.

\emph{Instruction count.} Dynamic counts use VLEN${}=128$ and LMUL${}=1$. To
cover all matrices we use a fast implementation with the same arithmetic as the
simulator and an analytical count; the test suite requires bit-identical
results and equal instruction counts between both on synthetic matrices,
including masked, overflowing, and multi-strip rows.

\emph{Ablation.} A variant of \code{vtdot.vv} that rounds every product and
partial sum to $2\mathrm{SEW}$, like \code{vfwmacc}, separates the effect of the
format from that of the exact accumulation.

\section{Results and Discussion}
\begin{figure}[t]
\centering
\begin{tikzpicture}
\begin{groupplot}[group style={group size=1 by 3, vertical sep=7mm},
  width=\columnwidth, height=3.3cm, xmin=0, xmax=100, grid=major,
  grid style={gray!25}, tick label style={font=\scriptsize},
  label style={font=\scriptsize}, title style={font=\scriptsize, yshift=-1.5ex},
  legend style={font=\scriptsize, at={(0.02,0.95)}, anchor=north west},
  unbounded coords=jump, no markers, every axis plot/.append style={thick}]
\nextgroupplot[title={8 bits}, ymin=-4, ymax=2, y filter/.expression={y<-4 ? -4 : y}]
\addplot table[x=share, y=T8, col sep=comma] {\BenchCDF}; \addlegendentry{T8}
\addplot table[x=share, y=E4M3, col sep=comma] {\BenchCDF}; \addlegendentry{E4M3}
\addplot table[x=share, y=E5M2, col sep=comma] {\BenchCDF}; \addlegendentry{E5M2}
\nextgroupplot[title={16 bits}, ymin=-8, ymax=2, y filter/.expression={y<-8 ? -8 : y}, ylabel={$\log_{10}$ relative error}]
\addplot table[x=share, y=T16, col sep=comma] {\BenchCDF}; \addlegendentry{T16}
\addplot table[x=share, y=FP16, col sep=comma] {\BenchCDF}; \addlegendentry{binary16}
\nextgroupplot[title={32 bits}, ymin=-14, ymax=2, y filter/.expression={y<-14 ? -14 : y}, xlabel={matrices (\%)}]
\addplot table[x=share, y=T32, col sep=comma] {\BenchCDF}; \addlegendentry{T32}
\addplot table[x=share, y=FP32, col sep=comma] {\BenchCDF}; \addlegendentry{binary32}
\end{groupplot}
\end{tikzpicture}
\caption{\BenchDraft Cumulative distribution of the SpMV relative error against
binary128, by format. A curve ends where the remaining matrices produced NaN or
an infinity; smaller errors, including exact results, are drawn at the
bottom of the axis.}
\label{fig:cdf}
\end{figure}

\begin{table}[t]
\caption{\BenchDraft SpMV results per format (\BenchMatrices{} matrices).
Median $\log_{10}$ of the relative error, share of matrices with error below
100\%, share with NaN or infinity, and vector instructions per nonzero}
\label{tab:results}
\centering
\setlength{\tabcolsep}{4pt}
\begin{tabular}{lrrrr}
\toprule
Format & Median & $<100\%$ & NaN/$\infty$ & Vec./nnz\\
\midrule
T8 & $\BenchMedianTEight$ & \BenchBelowOneTEight\% & \BenchInfTEight\% & \BenchVecPerNnzTEight\\
E4M3 & $\BenchMedianEFourMThree$ & \BenchBelowOneEFourMThree\% & \BenchInfEFourMThree\% & \BenchVecPerNnzEFourMThree\\
E5M2 & $\BenchMedianEFiveMTwo$ & \BenchBelowOneEFiveMTwo\% & \BenchInfEFiveMTwo\% & \BenchVecPerNnzEFiveMTwo\\
\midrule
T16 & $\BenchMedianTSixteen$ & \BenchBelowOneTSixteen\% & \BenchInfTSixteen\% & \BenchVecPerNnzTSixteen\\
binary16 & $\BenchMedianFPSixteen$ & \BenchBelowOneFPSixteen\% & \BenchInfFPSixteen\% & \BenchVecPerNnzFPSixteen\\
\midrule
T32 & $\BenchMedianTThirtyTwo$ & \BenchBelowOneTThirtyTwo\% & \BenchInfTThirtyTwo\% & \BenchVecPerNnzTThirtyTwo\\
binary32 & $\BenchMedianFPThirtyTwo$ & \BenchBelowOneFPThirtyTwo\% & \BenchInfFPThirtyTwo\% & \BenchVecPerNnzFPThirtyTwo\\
\bottomrule
\end{tabular}
\end{table}

\emph{Accuracy.} Fig.~\ref{fig:cdf} and Table~\ref{tab:results} summarize the
\BenchMatrices{} matrices. At 8 and 16 bits the decisive difference is dynamic
range. Every takum result was finite, and no matrix entry had to be saturated.
E4M3, whose largest value is 448, gave NaN for \BenchInfEFourMThree\% of the
matrices, E5M2 for \BenchInfEFiveMTwo\%, and binary16 for \BenchInfFPSixteen\%;
these formats also flushed nonzero entries to zero in \BenchFlushEFourMThree,
\BenchFlushEFiveMTwo, and \BenchFlushFPSixteen{} matrices. Where they stay in
range, however, the IEEE formats are usually more accurate. Compared matrix by
matrix, T8 has a smaller error than E4M3 in \BenchBetterTEightEFourMThree{} of
the \BenchInRangeTEightEFourMThree{} matrices where E4M3 is finite and a larger
one in \BenchWorseTEightEFourMThree; T16 beats binary16 in
\BenchBetterTSixteenFPSixteen{} of \BenchInRangeTSixteenFPSixteen{} and loses in
\BenchWorseTSixteenFPSixteen. This is the trade-off of tapered precision: takums
spend bits on range, and near unity they keep fewer fraction bits than an IEEE
format of the same width~\cite{takum}. T8 also has errors of 100\% or more in
\BenchAboveOneTEight\% of the matrices, caused by cancellation in rows whose
exact sum is small. At 32 bits both formats cover all the data and T32 is more
accurate: a smaller error in \BenchBetterTThirtyTwoFPThirtyTwo{} matrices, a
larger one in \BenchWorseTThirtyTwoFPThirtyTwo, and a median of
$\BenchMedianTThirtyTwo$ against $\BenchMedianFPThirtyTwo$, in line
with~\cite{simd}.

\emph{Exact accumulation.} Rounding every step to $2\mathrm{SEW}$ instead of
using \code{vtdot.vv} changed the result in only
\the\numexpr\BenchBetterTEightTEightStep+\BenchWorseTEightTEightStep\relax{} T8,
\the\numexpr\BenchBetterTSixteenTSixteenStep+\BenchWorseTSixteenTSixteenStep\relax{} T16,
and \the\numexpr\BenchBetterTThirtyTwoTThirtyTwoStep+\BenchWorseTThirtyTwoTThirtyTwoStep\relax{}
T32 matrices, and in T8 the exact sum was better in
\BenchBetterTEightTEightStep{} of them and worse in
\BenchWorseTEightTEightStep. With $x$ of ones every product is exact and most
rows are short, so fusion has little effect on SpMV accuracy here; its benefit
in this kernel is the instruction count.

\emph{Instruction count.} \code{vtdot.vv} replaces the multiply, the change of
\code{vtype}, and the reduction of the IEEE kernel, and the FP8 kernel also needs
two widening conversions per strip. Over all matrices the takum kernel executes
\BenchVecReductionEFourMThree\% fewer vector instructions than the FP8 kernels,
\BenchVecReductionFPSixteen\% fewer than binary16 and
\BenchVecReductionFPThirtyTwo\% fewer than binary32. Counting the scalar loop
control, which is identical for all formats, the total reductions are
\BenchTotalReductionEFourMThree\%, \BenchTotalReductionFPSixteen\%, and
\BenchTotalReductionFPThirtyTwo\%. Fewer instructions do not imply fewer cycles:
a fused dot product is a longer operation in hardware, which this functional
model does not measure.

\emph{Control state.} No kernel instruction reads \code{frm} or sets a flag in
the takum version. The IEEE kernels depend on the dynamic rounding mode, and
their overflow appears only as a NaN or an infinity in the data or as a sticky
flag that the code must test. Takum saturation keeps results finite but
silent; it did not occur in these matrices, and an architecture that needs to
detect it would have to keep one flag, which is still less state than
\code{fcsr}.

\emph{Limitations.} The model is functional and written in Python with exact
rationals, so it says nothing about speed or hardware cost. The FP8 baseline
uses conversions that ratified RVV lacks. Only $x$ of ones was evaluated, and
SpMV exercises one access pattern.

\section{Conclusion}
Takums fit into RVV as a format selected by SEW, with five instructions under a
custom opcode and no new control or status state. On \BenchMatrices{} real
matrices, 8- and 16-bit takums always gave finite results where OFP8 and
binary16 failed for a quarter to a third of the matrices, at the cost of lower
accuracy than those formats where they stay in range; at 32 bits takum was the
more accurate format. The fused widening dot product shortens the SpMV inner
loop compared with every IEEE kernel. Natural next steps are a cycle-level or RTL model of \code{vtdot.vv},
the logarithmic takum variant, and kernels beyond SpMV.

\section*{Reproducibility}
The repository contains the model, tests, scripts, and results.
\code{make verify} repeats the phase~1 checks and the test suite, \code{make
fetch} and \code{make bench} download the matrices and rerun the benchmark
(about three minutes on 16 cores), and
\code{make paper} builds this paper from the result files. Source:
\url{https://github.com/pedrohcali/aoc-takum-rvv}.

\begin{thebibliography}{10}
\bibitem{ieee754}
``IEEE Standard for Floating-Point Arithmetic,'' \emph{IEEE Std 754-2019},
Jul. 2019, doi: 10.1109/IEEESTD.2019.8766229.
\bibitem{ofp8}
P. Micikevicius, S. Oberman, P. Dubey \emph{et al.}, ``OCP 8-bit Floating Point
Specification (OFP8),'' Open Compute Project, Jun. 2023.
\bibitem{simd}
L. Hunhold, ``Streamlining SIMD ISA Extensions with Takum Arithmetic:
A Case Study on Intel AVX10.2,'' in \emph{Proc. 14th Int. Conf. Modern Circuits
and Systems Technologies (MOCAST)}, Jun. 2025, pp. 1 to 6.
\bibitem{takum}
L. Hunhold, ``Beating Posits at Their Own Game: Takum Arithmetic,'' in
\emph{Next Generation Arithmetic}, LNCS, vol. 14666. Springer, 2024,
pp. 1 to 51, doi: 10.1007/978-3-031-72709-2\_1.
\bibitem{rvv}
``RISC-V `V' Vector Extension,'' version 1.0, RISC-V International, Sep. 2021.
\bibitem{rvunpriv}
``The RISC-V Instruction Set Manual, Volume I: Unprivileged Architecture,''
RISC-V International, 2024.
\bibitem{posit}
J. L. Gustafson and I. Yonemoto, ``Beating Floating Point at Its Own Game:
Posit Arithmetic,'' \emph{Supercomputing Frontiers and Innovations}, vol. 4,
no. 2, pp. 71 to 86, 2017.
\bibitem{sparsesolvers}
L. Hunhold and J. Quinlan, ``Evaluation of Bfloat16, Posit, and Takum
Arithmetics in Sparse Linear Solvers,'' in \emph{Proc. IEEE ARITH}, 2025,
arXiv: 2412.20268.
\bibitem{takumvhdl}
L. Hunhold, ``Design and Implementation of a Takum Arithmetic Hardware Codec
in VHDL,'' 2024, arXiv: 2408.10594.
\bibitem{fphub}
J. Hormigo \emph{et al.}, ``FPHUB-RISCV: HUB Floating-Point Unit in RISC-V
Platform: Format definition,'' in \emph{RISC-V Summit Europe}, Paris, France,
May 2025.
\bibitem{libtakum}
L. Hunhold, ``libtakum: Takum arithmetic C99 reference implementation,''
version 2.0.0, revision \code{776ac0c}.
[Online]. Available: \url{https://github.com/takum-arithmetic/libtakum}
\bibitem{matrixmarket}
R. F. Boisvert, R. Pozo, K. Remington, R. F. Barrett, and J. J. Dongarra,
``Matrix Market: A Web Resource for Test Matrix Collections,'' in
\emph{Quality of Numerical Software}. Springer, 1997, pp. 125 to 137.
[Online]. Available: \url{https://math.nist.gov/MatrixMarket/}
\bibitem{suitesparse}
T. A. Davis and Y. Hu, ``The University of Florida Sparse Matrix Collection,''
\emph{ACM Trans. Math. Softw.}, vol. 38, no. 1, pp. 1 to 25, 2011.
\end{thebibliography}
\end{document}
